\documentclass[11pt,a4paper]{article}
\usepackage[utf8]{inputenc}
\usepackage[T1]{fontenc}
\usepackage{amsmath,amssymb,amsthm}
\usepackage{geometry}
\usepackage{enumitem}
\usepackage{fancyhdr}
\usepackage{titlesec}
\usepackage{hyperref}

\geometry{margin=1in}

\pagestyle{fancy}
\fancyhf{}
\fancyhead[C]{\small Cayley's Collected Mathematical Papers---Vol.~X, pp.~601--640}
\fancyfoot[C]{\thepage}
\renewcommand{\headrulewidth}{0.4pt}

\theoremstyle{definition}
\newtheorem{problem}{Problem}
\newtheorem*{solution}{Solution}
\newtheorem*{remark}{Remark}
\newtheorem*{note}{Note}

\titleformat{\section}{\large\bfseries}{\thesection.}{0.5em}{}
\titleformat{\subsection}{\normalsize\bfseries}{\thesubsection.}{0.5em}{}

\newcommand{\ditto}{""}

\begin{document}

\begin{center}
\Large\bfseries
PROBLEMS AND SOLUTIONS.
\end{center}

\thispagestyle{empty}

\vspace{1em}

\noindent\textbf{[705]}

\section*{Problems 4354--4458 [Vol.~XXII., July to December, 1874]}

\subsection*{4354. (Proposed by R.~Tucker, M.A.)\---Solve the equations}

\begin{align}
-x^2 + xy + xz &= a = -20, \\
-y^2 + xy + yz &= b = -8, \\
-z^2 + xz + yz &= c = -4.
\end{align}

\begin{note}[By Professor Cayley. Note on Question 4354.]
A question of simple algebra such as this, becomes more interesting when interpreted geometrically: thus, writing the equations in the form
\[
\begin{aligned}
-x^2 + xy + xz &= aw^2, \\
zx + zy - z^2 &= cw^2, \\
yx - y^2 + yz &= bw^2,
\end{aligned}
\]
and then putting for shortness
\[
\alpha = -a + b + c, \quad
\beta = a - b + c, \quad
\gamma = a + b - c,
\]
the solutions obtained are
\[
x : y : z : w = \alpha a : \beta b : \gamma c : (\alpha\beta\gamma)^{\frac{1}{2}},
\]
and
\[
x : y : z : w = \alpha a : \beta b : \gamma c : -(\alpha\beta\gamma)^{\frac{1}{2}};
\]
say these are
\[
[\alpha a, \; \beta b, \; \gamma c, \; (\alpha\beta\gamma)^{\frac{1}{2}}]
\]
and
\[
[\alpha a, \; \beta b, \; \gamma c, \; -(\alpha\beta\gamma)^{\frac{1}{2}}].
\]

But the equations are also satisfied by
\[
(x=0, \; y=z, \; w=0), \quad\text{or}\quad (y=0, \; z=x, \; w=0), \quad\text{or}\quad (z=0, \; x=y, \; w=0),
\]
that is, $(0,1,1,0)$, $(1,0,1,0)$, $(1,1,0,0)$; and the three having in common the line $x=y=z$, $w=0$, these are points; the last mentioned two points each once, and the first twice: $2 + 3 \cdot 2 = 8$.

The three equations represent quadric surfaces, each two of them intersecting in a proper curve; the three have in common the line $x=y=z$, $w=0$. To verify this, observe that, at each of the three points, the surfaces have a common tangent line; this is made up of the tangent planes of each two of the surfaces at the point, and the curve of intersection of any two of the surfaces therefore touches the third surface; wherefore the point counts for two intersections.

In fact, taking $(X,Y,Z,W)$ as current coordinates, the equations of the tangent planes at $(x,y,z,w)$ are
\[
\begin{aligned}
&X(2x-y-z) - Y\cdot x - Z\cdot x + 2aWw = 0, \\
&-X\cdot y + Y(-x+2y-z) - Z\cdot y + 2bWw = 0, \\
&-X\cdot z - Y\cdot z + Z(-x-y+2z) + 2cWw = 0;
\end{aligned}
\]
hence at the point $(0,1,1,0)$ these equations are
\[
-X + Y - Z = 0, \quad -X - Y + Z = 0, \quad -2Z = 0,
\]
which three planes meet in the line $X=0$, $Y-Z=0$; and similarly for the other two of the three points.
\end{note}

\newpage

\subsection*{4458. (Proposed by Professor Cayley.)\---[Vol.~XXII., pp.~60--64.]}

Find the intersections of the two quartic curves
\begin{align}
\label{eq:4458-1}
(\lambda ab - xy)^2 &= abx(a-y)(b-y), \\
\label{eq:4458-2}
(\mu ab - xy)^2 &= aby(a-x)(b-x);
\end{align}
the curves are the trace of the surfaces in some particular cases; for instance, when $a=1$, $b=2$, $\lambda=\mu=-2$.

\begin{solution}[Solution by the Proposer.]
1.~The 16 intersections are made up as follows: 5 points at infinity on the line $x=0$, 5 at infinity on the line $y=0$, the two points $(x=a,\; y=b)$, $(x=b,\; y=a)$, and 4 other points; $16 = 5+5+2+4$. As to the points at infinity, observe that, as regards the first curve, this line is the same for the point at infinity on the line $x=0$, which is a flecnode, having a cusp, and the line $y=b$ for the tangent to the flecnodal branch; this point counts as $2+3=5$ intersections; and, similarly, as regards the second curve, the point at infinity on the line $y=0$.

It remains to find the coordinates of the 4 points of intersection. Assume $xy = ab\omega$, then the equations become
\[
\lambda(1-\omega)\cdot ab = ax + y - (a+b)x, \quad
\mu(1-\omega)\cdot ab = by + x - (a+b)y,
\]
hence, eliminating successively $x$ and $y$, the factor $1-\omega$ divides out, this for $(x=a, y=b)$ for which obviously $\omega = 1$; and the equations become
\[
(\lambda - \lambda\omega)(1-\omega) + (a+b)\omega = (1+\omega)x, \quad
(\mu - \mu\omega)(1-\omega) + (a+b)\omega = (1+\omega)y.
\]
Multiplying these two equations together, and substituting for $xy$ its value $ab\omega$, we find
\[
\{2\lambda\mu(1-\omega) + (a+b)\}(1-\omega)^2 + (a+b)^2\omega^2 - (1+\omega)^2 ab = 0.
\]
Write, for shortness,
\[
p = (\lambda + \mu)(a+b) - \lambda\mu,
\]
then, dividing by $\omega$ and writing $\omega + \frac{1}{\omega} = \theta$, the equation is
\[
(2\lambda\mu + p)(\theta - 2) + (a+b)^2 - ab(\theta + 2) = 0;
\]
viz. this is a quadric equation for $\theta$. But, instead of this equation thus for $\theta$, it is convenient to introduce $\epsilon = \frac{\theta-1}{\theta+1}$, or
\[
\theta = \frac{1+\epsilon}{1-\epsilon};
\]
the equation becomes
\[
\{2\lambda\mu(1-\epsilon) + p(1+\epsilon)\}(1-\epsilon) + (a+b)^2(1-\epsilon)^2 - 4ab(1-\epsilon) = 0,
\]
or
\[
\theta^2\{(a+b)^2 - 4(p - 2\lambda\mu)\} + \theta\{-2a^2 - 2b^2 + 4(p + 2\lambda\mu)\} + (a-b)^2 = 0;
\]
viz. substituting for $p$ its value, this is
\[
\theta^2(a+b-2\lambda-2\mu)^2 + 2\theta\{-a^2-b^2+2(\lambda+\mu)(a+b)-2(\lambda-\mu)^2\} + (a-b)^2 = 0;
\]
or if we write
\[
A = a^2 - 2a(\lambda+\mu) + (\lambda-\mu)^2, \quad
B = b^2 - 2b(\lambda+\mu) + (\lambda-\mu)^2,
\]
this is
\[
\theta^2(a+b-2\lambda-2\mu)^2 - 2(A+B)\theta + (a-b)^2 = 0,
\]
whence
\[
\{(a-b)^2 - (A+B)\theta\}^2 = \theta^2\{(A+B)^2 - (a-b)^2(a+b-2\lambda-2\mu)^2\}
= \theta^2\{(A+B)^2 - (A-B)^2\} = 4AB\theta^2;
\]
viz. taking for convenience the sign $-$ on the right-hand side, this is
\[
(a-b)^2 - (A+B)\theta = -2\theta\sqrt{AB};
\]
and we have thus
\[
\theta = \frac{(a-b)^2}{(\sqrt{A} - \sqrt{B})^2},
\]
that is,
\[
\frac{\theta-1}{\theta+1} = \frac{1}{a-b}\cdot\frac{A-B}{\sqrt{A}+\sqrt{B}} = \frac{\sqrt{A}-\sqrt{B}+a-b}{a-b}.
\]

We may write
\[
x = \frac{1}{\omega}\{\epsilon(\omega-1) + \tfrac{1}{2}(a+b) + \tfrac{1}{2}(a+b-2\lambda-2\mu)\theta\},
\]
\[
y = \frac{1}{\omega}\{\epsilon(\omega-1) + \tfrac{1}{2}(a+b) + \tfrac{1}{2}(a+b-2\lambda-2\mu)\theta\};
\]
whence also $x-y = (\mu-\lambda)(\omega-1)$, as is seen at once from the original equations; then we have
\[
\tfrac{1}{2}(a-b)\{\tfrac{1}{2}(a+b-2\lambda-2\mu)\}\theta = \tfrac{1}{2}(\sqrt{A}-\sqrt{B})(\sqrt{A}+\sqrt{B}-\frac{a-b}{2\epsilon})
\]
and the values are
\[
x = \frac{(a-b)(\mu-\lambda) + b\sqrt{A} - a\sqrt{B}}{2\lambda(a-b)}, \quad
y = \frac{(a-b)(\lambda-\mu) + b\sqrt{A} - a\sqrt{B}}{2\mu(a-b)},
\]
which may be expressed in the more simple form
\[
x = \frac{1}{\lambda}(a+\lambda-\mu+\sqrt{A})(b+\lambda-\mu+\sqrt{B}),
\]
\[
y = \frac{1}{\mu}(a-\lambda+\mu+\sqrt{A})(b-\lambda+\mu+\sqrt{B}),
\]
the transformations depending on the identity
\[
\sqrt{A}\sqrt{B} - a + b = ab - (\lambda+\mu)(a+b) + (\lambda-\mu)^2 + \sqrt{A}(b-\lambda-\mu) + \sqrt{B}(a-\lambda-\mu) + \sqrt{AB},
\]
which is easily verified. Of course, since the signs of $\sqrt{A}$, $\sqrt{B}$ are arbitrary, we have 4 systems of values of $(x,y)$, which is right.

In the original equations, for $a,b,\lambda,\mu,x,y$, write $k^{-2}, 1, \lambda x^2, -\mu y^2, x^2, -y^2$; then the equations become
\[
\lambda x^2(1+k^2y^2) = x^2(1+y^2)(1+k^2y^2), \quad
\mu y^2(1+k^2x^2) = y^2(1-x^2)(1-k^2x^2),
\]
and we thence have
\[
x^2 + y^2 = \frac{\lambda\sqrt{(1+y^2)(1+k^2y^2)} + \mu y\sqrt{(1-x^2)(1-k^2x^2)}}{1+k^2x^2y^2} = \operatorname{sn}(\alpha+\beta i);
\]
viz. if $x = \operatorname{sn}\alpha$ (sinam $\alpha$), $iy = \operatorname{sn}\beta$, this is $x^2+y^2 = \operatorname{sn}(\alpha+\beta i)$; viz. the problem is (for a given modulus $k$, less than 1, as usual, to be assumed positive, and the quantity $\lambda/\mu$ the form $\operatorname{sn}(\alpha+\beta i)$. The proper solution is that in which the signs of the radicals are each $+$; it may be shown that in this case the value of $x^2$ is positive and less than 1, viz. the value of $y^2$ is positive.

The values thus are
\[
x^2 = \tfrac{1}{2}(A'-B'-\sqrt{A'}\sqrt{B'}), \quad
y^2 = \tfrac{1}{2}(A'-B'+\sqrt{A'}\sqrt{B'}),
\]
where
\[
A' = 1 - 2\lambda^2 + 2\mu^2 + (\lambda^2+\mu^2)^2, \quad
B' = \tfrac{1}{2} - \tfrac{1}{2}\lambda^2 + \tfrac{1}{2}\mu^2 + (\lambda^2+\mu^2)^2.
\]

The solution is really equivalent to that given by Richelot (Crelle, t.~XLV., 1853, p.~225). To verify this partially, observe that, writing $a,r$ for Richelot's $\tan^2\phi$, $\tan^2\psi$, we have
\[
\frac{x^2}{y^2} = \frac{a}{r}; \quad
\lambda = \frac{1}{\sqrt{kr}}, \quad \mu = \frac{1}{\sqrt{ka}},
\]
giving
\[
\frac{1}{2}(\lambda+\mu)\sqrt{A'}, \quad
\frac{1}{2}(\lambda-\mu)\sqrt{B'};
\]
whence
\[
2\sigma\lambda = 1+\lambda^2+\mu^2-\sqrt{A'}, \quad
2\tau\mu = 1+\lambda^2+\mu^2-\sqrt{B'};
\]
or the above value of $x^2$ is $= \frac{\sigma}{\sigma+\tau}$, agreeing with his; the value of $y^2$ is, however, presented under a somewhat different form.

\medskip

\textit{2.}~The curves are
\[
(2-xy)^2 = 2x(1-y)(2-y) \quad\text{and}\quad (2-xy)^2 = y(1-x)(2-x),
\]
each passing through the points $(1,2)$, $(2,1)$; the four points of intersection found by the foregoing general theory are all real, viz.
\[
x = \tfrac{1}{4}(2+\sqrt{3})(5+\sqrt{17}) = \text{say } 17.00 \text{ and } 1.65,
\]
\[
y = -\tfrac{1}{4}(-1+\sqrt{3})(-1+\sqrt{17}) = \text{say } +0.94 \text{ and } +2.13,
\]
and
\[
x = \tfrac{1}{4}(2-\sqrt{3})(5-\sqrt{17}) = \text{say } -0.12 \text{ and } -3.49,
\]
\[
y = -\tfrac{1}{4}(-1-\sqrt{3})(-1-\sqrt{17}) = \text{say } +1.22 \text{ and } -0.57.
\]

The equation of the curve (1) may also be written in the forms
\[
x(y^2-2x) + 2yx - 4x + 4 = 0, \quad xy^2 + x(-2y^2+2y-4) + 4 = 0.
\]
The original form shows that, if $x$ is negative (but by a further examination it appears that between these limits of $y$, $y$ being outside these limits, then $x$ is positive; in fact, the whole curve lies on the positive side of the axis of $y$. And then the inspection of the first quadric equation shows that the lines $x=\infty$ and $x=2$ are each an asymptote. The point at infinity on the axis of $y$ is a flecnode, the tangent to the flecnodal branch being $x=0$, and that of the ordinary branch $x=2$. Similarly, from the second quadric equation, the point at infinity on the axis of $x$ is in fact a flecnode, having the line $y=0$ for the tangent to the flecnodal branch, and $y=2$ for that of the other branch. It further appears that the line $y=0$ is in fact a cusp, having the axis $x$ for the cuspidal tangent.

The equation of the curve (2) may also be written in the forms
\[
x^2y + (x^2-7x+2)y + 4 = 0, \quad y^2x + (y^2-7y+2)x + 4 = 0.
\]
The original form shows that, if $y$ is between 1 and 2, the points at $x=2$ and $\frac{1}{2}$, and the easy to see that the between these limits, $y=1$ and $y=2$, $x$ is negative; and if $x$ between 1 and 2, is positive; in fact, the existence of an oval, meeting the line $y=1$ and $y=2$ in the points $x=2$ and $x=\frac{1}{2}(25+\sqrt{113}) = \text{say } 2.2$ and $0.9$. The remainder of the curve lies wholly below the line $y=0$. The first quadric equation shows the asymptote $y=0$; the point at infinity on the axis of $x$ is a cusp, having the axis $y$ for the cuspidal tangent. The second quadric equation shows the asymptotes $y=0$, $y=-1$; the point at infinity on the axis of $x$ is in fact a flecnode, having the line $y=0$ for the tangent to the flecnodal branch, and $y=1$ for that of the other branch. It is further seen that there are two vertical tangents $x = \frac{1}{2}(11 \pm \sqrt{113}) = 10.8$ or $0.2$; the former of these touches a branch between the two asymptotes $y=0$, $y=-1$; the latter one of the branches belonging to the cuspidal asymptote $x=2$; this last branch cuts the asymptote $y=0$ at $y=-2$, and then, cutting the asymptote $y=-1$ and $x=-\frac{1}{3} (= -0.3)$, goes on to touch at infinity the asymptote $y=0$. It is now easy to trace the curve.

The figure shows the two curves. The curve (1) is shown by a continuous line, the curve (2) by a thick dotted line; the points 1, 2, 3, 4 show the above mentioned four intersections of the curves; the point 1 and the dotted branch through it are drawn considerably out of their true positions; viz. as above appearing, the $x$-coordinate of 1 is $=17.00$, and the equation of the vertical tangent to the branch is $x=10.8$.
\end{solution}

\newpage

\subsection*{4620. (Proposed by A.~B.~Evans, M.A.)\---[Vol.~XXII., pp.~78, 79.]}

Find the least integral values of $x$ and $y$ that will satisfy the equation $x^2 - 953y^2 = -1$.

\begin{solution}[Solution by Professor Cayley.]
The values are given in Degen's Tables, viz.
\[
x = 2746864744198985747130, \quad y = 88979677\,488830275367615376.
\]
The work referred to is entitled ``Canon Pellianus, sive Tabula simplicissimam aequationis $y' = ax^2+1$ solutionem pro singulis numeri dati valoribus ab 1 usque ad 1000 in numeris rationalibus iisdemque integris exhibens. Auctore C.~F.~Degen, Hafniae (Copenhagen), 1817.''

Table I., pp. 3---106 gives, for all numbers 1 to 1000, the denominators, ($\eta$) the quotients of the convergent fraction of $\sqrt{a}$, and also the least values of $x$, $y$ which will satisfy the equation $x^2 - ay^2 = +1$. Thus
\[
a=953, \quad x = 2746864744198985747130, \quad y = 88979677\,488830275367615376.
\]
Table II., pp. 109---112, is described as giving for all those values of $a$ between 1 and 1000, for which there exists a solution of the equation $x^2 - ay^2 = -1$, the least values of $x$ and $y$ which will satisfy this equation; thus $a=953$, $x$ and $y$ as above. It is, however, to be noticed that the values $a = \beta^2+1$, $x=\beta$, $y=1$, for which there is the obvious solution, are omitted from the table. The reason for this appears, but the heading should have been different.
\end{solution}

\newpage

\section*{Problems 4528--4793 [Vol.~XXIII.--XXIV., 1875--1876]}

\subsection*{4528. (Proposed by Professor Cayley.)\---[Vol.~XXIII., pp.~18, 19.]}

A lottery is arranged as follows: there are $n$ tickets representing $a,b,c,\ldots$ pounds respectively. A person draws once, looks at his ticket; and if he pleases, draws again (out of the remaining $n-1$ tickets), looks at his ticket, and if he pleases draws again (out of the remaining $n-2$ tickets); and so on, drawing in all not more than $k$ times; and he receives the value of the last drawn ticket. Supposing that he regulates his drawings in the manner most advantageous to him according to the theory of probabilities, what is the value of his expectation?

\begin{solution}[Solution by the Proposer.]
Let the expression ``$a$ or $a$'' signify ``$a$ or $a$, whichever of the two is greatest,'' and let $M_1(a,b,c,\ldots)$ denote the mean of the quantities $(a,b,c,\ldots)$, viz. their sum divided by the number of them.

To fix the ideas, consider five quantities $a,b,c,d,e$, and write
\[
\begin{aligned}
M_1(a,b,c,d,e) &= M_1(a,b,c,d,e), \\
M_2(a,b,c,d,e) &= M_1\{a \text{ or } M_1(b,c,d,e), \; b \text{ or } M_1(a,c,d,e), \; \ldots, \; e \text{ or } M_1(a,b,c,d)\}, \\
M_3(a,b,c,d,e) &= M_1\{a \text{ or } M_2(b,c,d,e), \; b \text{ or } M_2(a,c,d,e), \; \ldots, \; e \text{ or } M_2(a,b,c,d)\},
\end{aligned}
\]
and so on. And the like in the case of any number of quantities $a,b,c,\ldots$.

Then the value of the expectation is $= M_k(a,b,c,\ldots)$.

For, when $k=1$, the value is obviously $=M_1(a,b,c,\ldots)$.

When $k=2$, if $a$ is drawn, the adventurer will be satisfied or he will draw again, according as $a$ or $M_1(b,c,\ldots)$ is greatest; viz. in this case the value of the expectation is ``$a$ or $M_1(b,c,\ldots)$.'' So if $b$ is drawn, the adventurer will be satisfied or he will draw again, as $b$ or $M_1(a,c,\ldots)$ is greatest; viz. ``$b$ or $M_1(a,c,\ldots)$''; and so on: and the value of the total expectation, the several cases being equally probable, is
\[
M_1\{a \text{ or } M_1(b,c,\ldots), \; b \text{ or } M_1(a,c,\ldots), \; \ldots\} = M_2(a,b,c,\ldots):
\]
and the like for $k=3$, $k=4$, \&c.

For instance, $a,b,c,d = 1,2,3,4$,
\[
\begin{aligned}
M_1(1,2,3,4) &= \tfrac{5}{2}, \\
M_2(1,2,3,4) &= M_1(1\text{ or }\tfrac{9}{2},\; 2\text{ or }\tfrac{7}{2},\; 3\text{ or }\tfrac{5}{2},\; 4\text{ or }\tfrac{3}{2}) = M_1(\tfrac{9}{2}, \tfrac{7}{2}, 3, 4) = \tfrac{15}{4}, \\
M_3(1,2,3,4) &= M_1(1\text{ or }4,\; 2\text{ or }\tfrac{15}{4},\; 3\text{ or }\tfrac{7}{2},\; 4\text{ or }\tfrac{5}{2}) = M_1(4, \tfrac{15}{4}, \tfrac{7}{2}, 4) = \tfrac{49}{16}, \\
M_4(1,2,3,4) &= M_1(1\text{ or }\tfrac{49}{16},\; 2\text{ or }4,\; 3\text{ or }\tfrac{15}{4},\; 4\text{ or }\tfrac{7}{2}) = M_1(\tfrac{49}{16}, 4, \tfrac{15}{4}, 4) = \tfrac{177}{64}.
\end{aligned}
\]

Or finally $M_1, M_2, M_3, M_4 = \tfrac{5}{2}, \tfrac{15}{4}, \tfrac{49}{16}, \tfrac{177}{64} = M_1, M_2, M_3, M_4$.

\medskip

\textbf{Cor.}---If the $a,b,c,\ldots$ denote penalties instead of prizes, then the solution is the same, except that ``$a$ or $a$'' must now denote ``$a$ or $a$, whichever of them is least.''
\end{solution}

\newpage

\subsection*{4581. (Proposed by the Rev.~M.~M.~U.~Wilkinson.)\---[Vol.~XXIII., pp.~47, 48.]}

A witness, whose statement is what he opines once in $m$ times, and whose opinion is correct once in $n$ times, asserted that the number of a note, issued by a bank universally known to have issued notes numbered from $B$ to $B+A-1$ inclusive, was $B+P$, where $P$ is either $0,1,2,\ldots,$ or $A-1$. Prove (1) that the probability that the note in question was that note is
\[
\frac{(m-1)(n-1)}{mn}\left(\frac{1}{A-1}\right),
\]
and (2) the probability that this last statement was correct.

\begin{remark}[Remark by Professor Cayley.]
There is a serious difficulty in this question. Try the answer in the case $m=10$, $n=10$. The witness is what he opines once out of 10 times---he is an atrocious liar once out of 10 times, that is, wrongly 9 times out of 10; he is in what he thinks, 9 times out of 10; and he opines rightly once out of 10 times. The witness says that the note was signed by $X$, and the chance of this being so is found to be $\frac{1+9}{10+9} = \frac{10}{19}$, or more than $\frac{1}{2}$; the larger part $\frac{9}{19}$ of this obtained as follows: the witness having said that the note was signed by $X$, the chances are 9 out of 10 that he opines, and, thinking the reverse, the chance is 9 out of 10 that he thought wrongly, viz. that the note was signed by $X$. But can the statement of such a witness create any probability in favour of the event?

The fallacy seems to consist in the assumption that the opinion is of that nature that the note was a definite number forming a series. But if the opinion be that the note was not 99, such opinions, the witness of 10 times out of 10 opinion true it is quite conceivable; in 9 times out of no amount of ingenuity would be able to know the number of the note, from the mere fact that the witness opines that the note is not a named number.

Suppose there are 500 notes, and that $n$ can have a determinate value irrespective of the event. I think that in calling the probability $= p_1p_2$, we mean the probability of A's making the statement, ``B says that the event did happen'' to be regarded as a simple statement, and the probability of the truth of the statement to be taken to be $= p_1p_2$. Mr Martin seems to introduce into his solution a hypothesis contradictory to the assumptions of the question.

I remark further that in my solution I assume that the event is of such a nature that, when there is any testimony in regard to it, the probability is determined by that testimony, irrespectively of the antecedent probability. This is quite consistent with the antecedent probability being, not zero, but $k$ (as it may very well be) indefinitely small; so if $k$ were $=0$, the whole probability $= p_1p_2$. But there is no absolutely reason for assigning any determinate value for $k$; and so the solutions which assume respectively $\beta=1$ and $\beta=0$, seem to me on this ground erroneous.
\end{remark}

\newpage

\subsection*{4638. (Proposed by Professor Cayley.)\---[Vol.~XXIII., p.~58.]}

Find the equation of the surface which is the envelope of the quadric surface $ax^2+by^2+cz^2+dw^2=0$, where $a,b,c,d$ are variable parameters connected by the equation $Abc+Bca+Cab+Fad+Gbd+Hcd=0$; and in the particular case in which the constants $A,B,C,F,G,H$ satisfy the condition $(AF)^{\frac{1}{2}}+(BG)^{\frac{1}{2}}+(CH)^{\frac{1}{2}}=0$.

\subsection*{4694. (Proposed by Professor Cayley.)\---[Vol.~XXIV., July to December, 1875, p.~41.]}

Taking $F$, $F'$ a pair of reciprocal points with respect to a circle, centre $O$; then if $F$, $F'$ are centres of force, each force varying as $(\text{distance})^{-2m}$, prove that (1) the resultant force upon any point $P$ on the circle is in the direction of a fixed point $S$ on the axis $OFF'$; and if, moreover, the forces at the unit of distance are as $(OF)^{2m-2}$ to $(OF')^{2m-2}$, then (2) the resultant force is proportional to $(SP)^{-2m} \cdot (PV)^{-2m+2}$, where $PV$ is the chord through $S$.

\newpage

\subsection*{4793. (Proposed by Professor Wolstenholme, M.A.)\---[Vol.~XXIV., pp.~72--74.]}

If $y = x^n (\log x)^r$, where $n$ and $r$ are integers, prove that
\[
\frac{d^{n+r}y}{dx^{n+r}} = \frac{r(r-1)\cdots(r-n+1)}{x^r} \cdot y
- \frac{r(r-1)(r-2)\cdots(r-n+2)}{x^{r-1}} \cdot \frac{dy}{dx}
+ \cdots
\]
the coefficients being
\[
\frac{r(r-1)\cdots(r-s+1)}{s!} \cdot \frac{(n+r)!}{(n+s)!}
\]
for the $s$-th term, so that the result may be symbolically written
\[
\frac{d^{n+r}y}{dx^{n+r}} = \left\{ \frac{d}{dx} \cdot \left( x\frac{d}{dx} - n \right) \right\}^r y,
\]
where $D$ denotes $\frac{d}{dx}$ and operates on $\frac{d^{n+r}y}{dx^{n+r}}$ only, and $\Delta$ operates on all terms after the $r$-th, vanishing since $\Delta^r x^n = 0$ when $r$ is an integer $>n$.

\begin{solution}[Solution by Professor Cayley.]
Since $y = x^n (\log x)^r$, therefore $(x\frac{d}{dx} - n)y = rx^n(\log x)^{r-1}$; repeating the same operation, we have $(x\frac{d}{dx} - n)^r y = [r]^r x^n$, whence $\frac{d^n}{dx^n}(x\frac{d}{dx} - n)^r y = [r]^r [n]^r$.

Now, for any value whatever of the function $y$, we have
\[
\frac{d^n}{dx^n}(x\frac{d}{dx} - n)^r y = Ax^n \frac{d^{n+r}y}{dx^{n+r}} + Bx^{n-1}\frac{d^{n+r-1}y}{dx^{n+r-1}} + Cx^{n-2}\frac{d^{n+r-2}y}{dx^{n+r-2}} + \&\text{c.},
\]
the coefficients $A,B,C,\ldots$ being functions, presumably of $r,n$, of the form of the function $y$. It will, however, appear that $A,B,C,\ldots$ are, in fact, functions of $r$ only.

To see how this is, observe that $(x\frac{d}{dx} - n)^r$ consists of a set of terms $(x\frac{d}{dx})^s$, ($s=0$ to $r$), where $(x\frac{d}{dx})^s$ denotes $s$ repetitions of the operation $x\frac{d}{dx}$; by a well-known theorem, this is $= [x\frac{d}{dx} + s - 1]^s$, where, after expansion of the factorial, $(x\frac{d}{dx})^t$ is to be replaced by $x^t\frac{d^t}{dx^t}$; thus
\[
(x\frac{d}{dx})^2 = (x\frac{d}{dx}+1)^2 = x^2\frac{d^2}{dx^2} + x\frac{d}{dx}, \quad
(x\frac{d}{dx})^3 = [x\frac{d}{dx}+2]^3 = x^3\frac{d^3}{dx^3} + 3x^2\frac{d^2}{dx^2} + 2x\frac{d}{dx}, \; \&\text{c.}
\]

Thus $(x\frac{d}{dx} - n)^r$ consists of a series of terms $x^t\frac{d^{t+s}}{dx^{t+s}}$, where $t+s=r$, or $t=r-s$; and the term is $x^{r-s}\frac{d^r}{dx^r}$, or $x^{r-s}\frac{d^r}{dx^r}$. Observing that $\frac{d^n}{dx^n}(x\frac{d}{dx} - n)^r$ consists of a series of terms of the form $\frac{d^n}{dx^n}x^{r-s}\frac{d^r}{dx^r}$, where the unaccented symbol operates on the $x$ and the accented symbol on the $y$; and putting $x^{r-s}\frac{d^r}{dx^r} = x^{r-s}\frac{d^{r-s}}{dx^{r-s}} \cdot \frac{d^s}{dx^s}$, that is, $x^{r-s}\frac{d^r}{dx^r} = x^{r-s}\frac{d^{r-s}}{dx^{r-s}} \cdot \frac{d^s}{dx^s}$; so that the term is $x^{r-s}\frac{d^{n+r-s}y}{dx^{n+r-s}}$.

To understand how the coefficients can be independent of $n$, take the particular case $r=2$; then we have here
\[
\frac{d^n}{dx^n}(x\frac{d}{dx} - n)^2 y = Ax^n\frac{d^{n+2}y}{dx^{n+2}} + Bx^{n+1}\frac{d^{n+1}y}{dx^{n+1}} + Cx^n\frac{d^ny}{dx^n}.
\]
The right-hand side is
\[
\frac{d^n}{dx^n}\{x^2\frac{d^2}{dx^2} - (2n-1)x\frac{d}{dx} + n^2\}y
= \{x^2\frac{d^{n+2}}{dx^{n+2}} + 2nx\frac{d^{n+1}}{dx^{n+1}} + (n^2-n)\frac{d^n}{dx^n} - (2n-1)x\frac{d^{n+1}}{dx^{n+1}} - n(2n-1)\frac{d^n}{dx^n} + n^2\frac{d^n}{dx^n}\}y;
\]
hence
\[
A=1, \quad B = 2n-(2n-1)=1, \quad C = (n^2-n)-n(2n-1)+n^2 = 0;
\]
and we thus see also how in this particular case the last coefficient is $=0$, viz. that we have
\[
\frac{d^n}{dx^n}(x\frac{d}{dx} - n)^2 y = x^2\frac{d^{n+2}y}{dx^{n+2}} + x\frac{d^{n+1}y}{dx^{n+1}},
\]
without any term in $\frac{d^ny}{dx^n}$.

To find the coefficients $A,B,C,\ldots$ generally, write $y=x^{n+r+\theta}$, then $x\frac{d}{dx} - n = r + \theta$, and consequently
\[
\frac{d^n}{dx^n}(x\frac{d}{dx} - n)^r y = (r+\theta)^r x^{n+\theta} = (r+\theta)^r[r+\theta+n]^n x^{n+\theta};
\]
whence
\[
(r+\theta)^r[r+\theta+n]^n = A[r+\theta+n]^{n+r} + B[r+\theta+n]^{n+r-1} + \cdots
\]
or, what is the same thing,
\[
(r+\theta)^r = A[r+\theta]^r + B[r+\theta]^{r-1} + \cdots
\]
Since every term on the left-hand side, and every term on the right-hand side, contains the factor $r+\theta$; it is not on the right-hand side any term $[r+\theta]^r$; dividing the equation by $r+\theta$, then becomes
\[
(r+\theta)^{r-1} = A[r+\theta-1]^{r-1} + B[r+\theta-1]^{r-2} + \cdots
\]
and we thus have the series
\[
Ux = (1+x)^r, \quad\text{where } r+\theta = 1+x,
\]
taking the terms in the reverse order,
\[
Ux = u_0 + \frac{x}{1}\Delta u_0 + \frac{x(x-1)}{1\cdot 2}\Delta^2 u_0 + \&\text{c.},
\]
the well-known one.

Hence, in general,
\[
\frac{d^{n+r}y}{dx^{n+r}} = \left\{ \frac{d}{dx}\left(x\frac{d}{dx} - n\right) \right\}^r y,
\]
where observe that the last term is $= x\frac{d^{n+1}y}{dx^{n+1}}$.

For the function $y=x^n(\log x)^r$, the value of each side is $=[r]^r[n]^r$.
\end{solution}

\newpage

\subsection*{4752. (Proposed by Professor Cayley.)\---[Vol.~XXIV., pp.~89--91.]}

Mr~Wolstenholme's Question 3067 may evidently be stated as follows: If $(a,b,c)$ are the coordinates of a point on the cubic curve
\[
a^3 + b^3 + c^3 = (b+c)(c+a)(a+b),
\]
and if
\[
x : y : z = a(-a+b+c) : b(a-b+c) : c(a+b-c);
\]
then $(x,y,z)$ are the coordinates of a point on the same cubic curve. This being so, it is required to find the geometrical relation of the two points to each other.

\begin{solution}[Solution by Professor Cayley.]
1.~On referring to Wolstenholme's Question 3067 (Reprint, Vol.~XIII., p.~70), it appears that the coordinates $(x,y,z)$ of the point may be expressed in the more simple form
\[
x : y : z = a(-a+b+c) : b(a-b+c) : c(a+b-c);
\]
viz. the given relation between $(a,b,c)$ being equivalent to
\[
a^3 + b^3 + c^3 - (b+c)(c+a)(a+b) = 0,
\]
we have
\[
\frac{x}{a(-a+b+c)} = \frac{y}{b(a-b+c)} = \frac{z}{c(a+b-c)};
\]
and thence
\[
\frac{-4abc}{x^2/a - (y^2/b + z^2/c)} = \cdots;
\]
and consequently
\[
x : y : z = x'(-x'+y'+z') : y'(x'-y'+z') : z'(x'+y'-z'),
\]
whence the transformation in question.

2.~In place of $(x,y,z)$ for the coordinates of the second point, writing $(x',y',z')$, and $(x,y,z)$ for those of the first point, the two points are connected by the relation
\[
x' : y' : z' = x(-x+y+z) : y(x-y+z) : z(x+y-z);
\]
and we thence at once deduce the converse relation
\[
x : y : z = x'(-x'+y'+z') : y'(x'-y'+z') : z'(x'+y'-z').
\]
Hence, writing
\[
\begin{aligned}
(-x+y+z, \; x-y+z, \; x+y-z) &= (\xi, \eta, \zeta), \\
(-x'+y'+z', \; x'-y'+z', \; x'+y'-z') &= (\xi', \eta', \zeta'),
\end{aligned}
\]
and similarly
\[
x' : y' : z' = x\xi' : y\eta' : z\zeta' = x^2\xi : y^2\eta : z^2\zeta;
\]
so that $\xi/\xi' = \eta/\eta' = \zeta/\zeta'$; the coordinates of the two points, we see that
\[
x : y : z = x'\xi' : y'\eta' : z'\zeta';
\]
regarding these as inverse points one of the other in regard to the triangle $\xi=0$, $\eta=0$, $\zeta=0$.

To complete the solution, we must introduce these new coordinates into the equation of the cubic curve. Writing this under the form
\[
Hxyz + 2(-x+y+z)(x-y+z)(x+y-z) = 0,
\]
and observing that the equation is $(\xi+\eta)(\eta+\zeta)(\zeta+\xi) + 2\xi\eta\zeta = 0$; viz. this is a cubic curve inverting into itself. And the two points in question are thus any two inverse points on this cubic curve.

3.~In regard to the original form, that the point $(x,y,z)$ defined by the equations
\[
\frac{x}{-a^2+b^2+c^2} = \frac{y}{a^2-b^2+c^2} = \frac{z}{a^2+b^2-c^2}
\]
lies on the cubic curve
\[
a^3+b^3+c^3-(b+c)(c+a)(a+b)=0,
\]
Professor Sylvester proceeds as follows:---Writing
\[
(x,y,z) = (a^4-(b^2-c^2)^2, \; b^4-(c^2-a^2)^2, \; c^4-(a^2-b^2)^2) = (A,B,C), \quad\text{suppose;}
\]
and
\[
F(a,b,c) = a^3+b^3+c^3-(b+c)(c+a)(a+b),
\]
he observes that the truth of the theorem depends on the identity
\[
F(A,B,C) + F(a,b,c) \cdot F(a,-b,c) \cdot F(a,b,-c) \cdot F(a,-b,-c) = 0,
\]
and that, in order to prove the identity generally, it is sufficient to prove it for the three cases $a^2=0$, $a^2=b^2+c^2$, $a^2=b^2$, which may be effected without difficulty.

But, for a general proof of the identity, write
\[
\lambda = b^2+c^2, \quad \mu = b^2-c^2,
\]
so that
\[
A = a^4-\mu^2, \quad B = (a^2+\mu)(-a^2+\lambda), \quad C = (-a^2+\lambda)(a^2-\mu).
\]
Whence
\[
\begin{aligned}
-F(A,B,C) &= -(a^4-\mu^2)^3 + 2(a^2\lambda\mu)^2(a^2+3a^2\mu-8a^2b^2c^2)(a^4-\mu^2)(a^2-\lambda) \\
&= a^{12} - 6\lambda a^{10} + (6\lambda^2+9\mu^2-8b^2c^2)a^8 + \lambda(-2\lambda^2-18\mu^2+8b^2c^2)a^6 \\
&\quad + \mu^2(18\lambda^2-3\mu^2+8b^2c^2)a^4 + \lambda\mu^2(-6\lambda^2-8b^2c^2)a^2 + \mu^6.
\end{aligned}
\]
Moreover
\[
F(a,b,c) = a\{a^2-(b+c)^2\} - (b+c)\{a^2-(b-c)^2\},
\]
therefore
\[
F(a,-b,-c) = a\{a^2-(b+c)^2\} + (b+c)\{a^2-(b-c)^2\};
\]
whence
\[
\begin{aligned}
F(a,b,c) \cdot F(a,-b,-c) &= a^2\{a^2-(b+c)^2\}^2 - (b+c)^2\{a^2-(b-c)^2\}^2 \\
&= a^4 - 3\gamma^2 a^2 + \gamma^2(\gamma^2+2\delta^2),
\end{aligned}
\]
if $\gamma = b+c$, $\delta = b-c$. By changing the sign of $c$, we interchange $\gamma$ and $\delta$, and we thus have
\[
F(a,b,-c) \cdot F(a,-b,c) = a^4 - 3\delta^2 a^2 + \delta^2(2\gamma^2+\delta^2);
\]
and the identity to be verified is thus
\[
[a^4 - 3\gamma^2 a^2 + \gamma^2(\gamma^2+2\delta^2)] \cdot [a^4 - 3\delta^2 a^2 + \delta^2(2\gamma^2+\delta^2)]
\]
\[
= a^{12} - 6\lambda a^{10} + \cdots
\]
the values of $\lambda$, $\mu$ in terms of $\gamma$, $\delta$ being
\[
\lambda = \tfrac{1}{2}(\gamma^2+\delta^2), \quad \mu = \gamma\delta;
\]
substituting these values on the right-hand side, the verification can be completed without difficulty.
\end{solution}

\newpage

\section*{Problems 4946--5208 [Vol.~XXV.--XXVII., 1876--1877]}

\subsection*{4946. (Proposed by Professor Cayley.)\---[Vol.~XXV., January to July, 1876, p.~82.]}

Show that the attraction of an indefinitely thin double convex lens on a point at the centre of one of its faces is equal to that of the infinite plate included between the tangent plane at the point and the parallel tangent plane of the other face of the lens.

\subsection*{5020. (Proposed by W.~S.~B.~Woolhouse, F.R.A.S.)\---[Vol.~XXVI., July to December, 1876, pp.~41, 42.]}

Let $v=\frac{1}{2}(n+1)$, $v'=\frac{1}{2}n(n-1)$, $v''=\frac{1}{2}n(n-1)(n-2)$, \&c.\ be the maximum coefficient and the sums of the first, second, third, \ldots differences of the coefficients of the expansion of the binomial $(1+x)^n$ taken as the central or maximum coefficient, and also let $S_1, S_2, S_3, \ldots, S_n$ be the coefficients; then show that the algebraic function
\[
x^n - S_1 x^{n-1} + S_2 x^{n-2} - S_3 x^{n-3} + \&\text{c.}
\]
is divisible by $(x-1)^n$ without remainder; and that the sum of the numerical coefficients of the quotient is equal to $1 \cdot 3 \cdot 5 \cdots (2n-1)$.

[See Solution to Question 1894, Reprint, vol.~v., p.~113.]

\begin{solution}[Solution by Professor Cayley.]
Mr~Woolhouse's elegant theorem depends ultimately on the property of triangular numbers $\phi(-n)$, so that, writing down the series $\phi(n) = \frac{1}{2}(n^2+n)$, $\phi(n+1) = \cdots$ we have, in fact,
\[
\ldots, 10, 6, 3, 1, 0, 0, 1, 3, 6, 10, \ldots
\]
a single continuous series, obtained by giving to $n$ the different negative and positive integer values, zero included.

Thus a particular case is
\[
(1+x)^6 - 5(1+x)^3 + 9(1+x) - 5 = 1 \cdot 3 \cdot 5x^2 + \&\text{c.}\,x^3 + \ldots,
\]
where, on the left-hand side, the exponents are the triangular numbers $\frac{1}{2}n(n+1)$, and the coefficients, after the first, are the differences of the binomial coefficients of the power $2n$ (in the particular case, $n=3$); viz. the binomial coefficients being $1, 6, 15, 20, 15, 6, 1$, the differences taken as far as they are positive are $5, 9, 5$.

Expanding the several terms and writing down only the coefficients, we have a diagram:
\[
\begin{array}{rrrrrrr}
1, & 6, & 15, & 20, & 15, & 6, & 1, \\
-5 & 1, & 3, & 3, & 1, \\ 
+9 & 1, & 1, \\-5 & 1,
\end{array}
\]
The theorem in the particular case depends on the identities
\[
1-5+9-5 = 0, \quad 6-15+9=0, \quad 15-20+5=0,
\]
or writing, as above, $h,g,f,e,d,c,b,a$ to denote the triangular numbers
\[
h^2-5g^2+9f^2-5e^2=0, \quad h^3-5g^3+9f^3-5e^3=0,
\]
\[
h^4-5g^4+9f^4-5e^4=0, \quad h^5-5g^5+9f^5-5e^5=1\cdot2\cdot3\cdot1\cdot3\cdot5.
\]

These may be replaced by
\[
\begin{aligned}
h^2-5g^2+9f^2-5e^2 &= 0, \\
h^3-5g^3+9f^3-5e^3 &= 0, \\
h^4-5f^4+9f^4-5e^4 &= 0, \\
h^5-5g^5+9f^5-5e^5 &= 1\cdot2\cdot3\cdot1\cdot3\cdot5.
\end{aligned}
\]

Consider any one of these, for instance the third; the function on the left-hand is
\[
\tfrac{1}{24}h^4 - (6-1)\tfrac{1}{6}g^4 + (15-6)\tfrac{1}{2}f^4 - (20-15)e^4,
\]
or, introducing the values $a,b,c,d,e,f,g,h$ which are any successive triangular numbers, this is in fact
\[
\tfrac{1}{24}(\xi+6)^4 - 6\cdot\tfrac{1}{6}(\xi+5)^4 + 15\cdot\tfrac{1}{2}(\xi+4)^4 - 20(\xi+3)^4 + 15(\xi+2)^4 - 6(\xi+1)^4 + 1\cdot\xi^4 = 0,
\]
an immediate consequence of the well-known theorem
\[
(\xi+6)^m - 6(\xi+5)^m + 15(\xi+4)^m - 20(\xi+3)^m + 15(\xi+2)^m - 6(\xi+1)^m + \xi^m = 0,
\]
for any value of $m$ up to $m=5$, and $= 1\cdot2\cdot3\cdot4\cdot5\cdot6$ for $m=6$.

We have thus all the equations except the last; and as regards the last equation, observe that the equation to be verified is
\[
\tfrac{1}{6!}[\tfrac{1}{2}(e+6)(e+5)]^2 - 6[\tfrac{1}{2}(e+5)(e+4)]^2 + \cdots = 1\cdot3\cdot5;
\]
viz. this may be replaced by
\[
\tfrac{1}{2}(e+6)^2 - 6(e+5)^2 + \cdots = 2\mu \cdot 2\cdot3\cdot1\cdot3\cdot5 = 2\cdot4\cdot6\cdot1\cdot3\cdot5 = 1\cdot2\cdot3\cdot4\cdot5\cdot6,
\]
which is right.

It is clear that the proof, although worked out on a particular case, is perfectly general; and Mr~Woolhouse's theorem is thus proved.
\end{solution}

\newpage

\subsection*{5079. (Proposed by Professor Cayley.)\---[Vol.~XXVI., pp.~77, 78.]}

Show that the curve
\[
\{(\beta-\gamma)^2-\delta^2\}\{(x-\beta)^2+y^2\} + \{(\beta+\gamma)^2-\delta^2\}\{(x+\beta)^2+y^2\} = \{(1-\xi)^2+\eta^2\}\{(x-\gamma-\delta)^2+y^2\}
\]
where $t = (\sqrt{-1})$ as usual, $\xi = \frac{\beta^2+\gamma^2-\delta^2}{(\beta+\gamma)^2-\delta^2}$, $\eta = \frac{-2\gamma\delta}{(\beta+\gamma)^2-\delta^2}$, is a real bicircular quartic having the axial foci $\beta$, $-\beta$, $\gamma+\delta$, $\gamma-\delta$.

\begin{solution}[Solution by the Proposer.]
Consider the equation
\[
(1+m\sqrt{-1})\{(x-\beta)^2+y^2\} + q(1-m\sqrt{-1})\{(x+\beta)^2+y^2\} = (\lambda+\mu\sqrt{-1})\{(x-\gamma-\delta)^2+y^2\}.
\]
This is
\[
\begin{aligned}
&(1+m\sqrt{-1})\{x^2+y^2-\beta^2-2\beta x\sqrt{-1}\} \\
&\quad + q(1-m\sqrt{-1})\{x^2+y^2-\beta^2+2\beta x\sqrt{-1}\} \\
&\quad - (\lambda+\mu\sqrt{-1})\{x^2+y^2-\beta^2+\beta^2+\gamma^2-\delta^2-2\gamma x - 2(\beta-\gamma)\delta\sqrt{-1}\} = 0,
\end{aligned}
\]
where, putting the imaginary part equal to zero,
\[
\begin{aligned}
&m(1-q)(x^2+y^2-\beta^2) - 2\beta(1+q)\sqrt{-1} \\
&\quad - \mu\{x^2+y^2-\beta^2+(\beta^2+\gamma^2-\delta^2)-2\gamma x\} + 2\lambda(\beta-\gamma)\delta = 0,
\end{aligned}
\]
which will be true identically if
\[
\begin{aligned}
m(1-q) - \mu &= 0, \\
-2\beta(1+q) + \mu\gamma + \lambda\delta &= 0, \\
-\mu(\beta^2+\gamma^2-\delta^2) - 2\lambda\gamma\delta &= 0.
\end{aligned}
\]
The last gives
\[
\lambda = \varrho(\beta^2+\gamma^2-\delta^2), \quad \mu = -2\varrho\gamma\delta, \quad \varrho \text{ arbitrary};
\]
and then
\[
\begin{aligned}
&2\beta(1-q)\beta = 2\varrho\gamma\delta(\beta^2+\gamma^2-\delta^2) = 2\varrho\delta(\beta^2-\gamma^2-\delta^2), \\
&m(1-q) = -2\varrho\gamma\delta;
\end{aligned}
\]
so that, putting $\varrho\delta = (1-q)\xi$, or $\varrho = (1-q)\xi/\delta$, we have
\[
\begin{aligned}
1 &= \varrho(\beta^2+\gamma^2-\delta^2) = (1-q)\xi\frac{\beta^2+\gamma^2-\delta^2}{\delta}, \\
m &= -2\varrho\gamma = (1-q)\xi\frac{-2\gamma}{\delta}, \\
\mu &= (1-q)\xi(-2\gamma), \\
\lambda &= (1-q)\xi(\beta^2+\gamma^2-\delta^2).
\end{aligned}
\]
Therefore
\[
1 \pm m\sqrt{-1} = (1-q)\left\{\frac{\beta^2+\gamma^2-\delta^2}{\delta} \mp \frac{2\gamma\sqrt{-1}}{\delta}\right\},
\]
\[
\lambda \pm \mu\sqrt{-1} = (1-q)\xi\{[\beta^2+(\gamma+\delta)^2] \mp [\beta^2+(\gamma-\delta)^2]\sqrt{-1}\};
\]
and the equation is
\[
\{(\beta-\gamma)^2-\delta^2\}\{(x-\beta)^2+y^2\} + \{(\beta+\gamma)^2-\delta^2\}\{(x+\beta)^2+y^2\} = \{(1-\xi)^2+\eta^2\}\{(x-\gamma-\delta)^2+y^2\},
\]
which is a real curve having the axial foci $\beta$, $-\beta$, $\gamma+\delta$, $\gamma-\delta$; viz. $\beta$ being a focus, and the curve being real, it is clear that $-\beta$ is also a focus.
\end{solution}

\newpage

\subsection*{5130. (Proposed by Professor Cayley.)\---[Vol.~XXVII., January to June, 1877, p.~20.]}

Show that the envelope of a variable circle, having its centre on a given conic and cutting at right angles a given circle, is a bicircular quartic; which, when the given conic and the circle have double contact, becomes a pair of circles; and, by means of the last-mentioned particular case of the theorem, connect together the porisms arising out of the two problems:
\begin{enumerate}[label=(\roman*)]
\item Given two conics, to find a polygon of $n$ sides inscribed in the one and each side touching the circumscribed about the other.
\item Given two circles, to find a closed series of $n$ circles, each touching the two circles and the two adjacent circles of the series.
\end{enumerate}

\subsection*{5208. (Proposed by Professor Sylvester.)\---[Vol.~XXVII., pp.~81--83.]}

Let the magnitude of any ramification signify the number of its branches, and let its partial magnitudes in respect to any node signify the magnitudes of the ramifications which come together at that node. If at any node the largest magnitude exceeds by $k$ the sum of the other magnitudes, let the node be called \textit{superior by} $k$, or be said to be of \textit{superiority} $k$; but if no magnitude exceeds the sum of the other magnitudes, let the node be called \textit{subequal}. Then the theorem is, in any ramification, either there is one and only one subequal node; or else there are two and only two nodes each superior by unity, these two nodes being contiguous.

\begin{solution}[Solution by Professor Cayley.]
The proof consists in showing that (1) there cannot be more than one subequal node; (2) there cannot be more than two nodes each superior by unity; and if there is one such node, then there is, contiguous to it, another such node; (3) starting from a node which is superior by more than unity, there is always a contiguous node which is either of smaller superiority, or else subequal; for, these theorems holding good, we can, by (3), always arrive at a node which is either subequal or else superior by unity; in the former case, by (1), the subequal node thus arrived at is unique; in the latter case, by (2), we have, contiguous to the node arrived at, a second node superior by unity; and we have thus a unique pair of nodes each superior by unity.

I will prove only (3), as it is easy to see that the like process applies to the proof of (1) and (2).

Let the whole magnitude be $n$; and suppose at a node $P$ which is superior by $k$, the largest magnitude is $\alpha$, and that the other magnitudes are, say $\beta, \gamma, \delta$. We have
\[
\alpha = \beta + \gamma + \delta + k;
\]
and since $n = \alpha + \beta + \gamma + \delta$, we have thence $n = 2\alpha - k$, or
\[
\alpha = \tfrac{1}{2}(n+k), \quad \beta + \gamma + \delta = \tfrac{1}{2}(n-k);
\]
clearly $k$ is even or odd, according as $n$ is even or odd.

Suppose now that we pass from $P$, along the branch of magnitude $\alpha$, to a node $Q$; and let the magnitudes for $Q$ be $\alpha', \beta', \gamma', \delta', \epsilon$, where $\alpha'$ denotes the magnitude for the branch $QP$. We have $\alpha' = \beta + \gamma + \delta + 1$ for the contiguous ramification: the branch $QP$ and of the ramifications of magnitudes $\beta, \gamma, \delta$ which meet in $P$. We have thus
\[
\alpha' = \tfrac{1}{2}(n-k) + 1 = \tfrac{1}{2}n - \tfrac{1}{2}(k-2)
\]
and thence
\[
\beta' + \gamma' + \delta' + \epsilon' = \tfrac{1}{2}n + \tfrac{1}{2}(k-2).
\]
Supposing here that $k$ is greater than 1, viz. that it is $=$ or $> 2$, $k-2$ is positive; and if $\alpha'$ be the greatest magnitude belonging to the node $Q$, this is a subequal node. But it may be that $\alpha'$ is not the greatest magnitude; supposing then that the greatest magnitude is $\beta'$, we have
\[
\beta' = \tfrac{1}{2}n + \tfrac{1}{2}(k-2) - \gamma' - \delta' - \epsilon',
\]
\[
\alpha' + \gamma' + \delta' + \epsilon' = \tfrac{1}{2}n - \tfrac{1}{2}(k-2) + \gamma' + \delta' + \epsilon',
\]
and thence
\[
\beta' - (\alpha' + \gamma' + \delta' + \epsilon') = k-2 - 2(\gamma' + \delta' + \epsilon');
\]
viz. $k-2$ that is, if the greatest branch of the node $Q$ is either subequal, or else, being superior, the superiority is at most $= k-2$. That is, from the node $P$, of superiority $=$ or $> 2$, we pass along the branch of greatest magnitude to the contiguous node $Q$, this is either subequal, or else of superiority less than that of $P$; which is the foregoing proposition (3).

The subequal node, and the two nodes of superiority 1, may be termed respectively the \textit{centre} and the \textit{bicentre}; the theorem thus is, every ramification has either a centre or else a bicentre.

But the centre and the bicentre here considered, due (as remarked by Professor Sylvester) to M.~Camille Jordan, and which we here consider, not the magnitude, but the length of a ramification, such distance: viz. length being measured by the number of branches to be travelled over in order to reach the most distant terminal node. The ramification has either a centre or else a bicentre of distance: viz. if, disregarding the other ramifications, it is a node such that, for two or more of the branches which proceed from it, the lengths are equal; and the bicentre a pair of contiguous nodes such that, from the two nodes, there are from each of them two or more ramifications the lengths of which are equal to each other and superior to those of the other ramifications.

It is very noticeable how close the agreement is between the proofs for the two kinds of centre. Say, as regards distance, if at any node the length of the longest branch exceeds by $k$ the length of the next longest branch or branches, then the node is superior by $k$, or is of the superiority $k$; but, if there are two or more longest branches, then the node is subequal. And say next, in regard to magnitude, if at any node the largest magnitude exceeds by $k$ the sum of all the other magnitudes, the node is superior by $k$, or has a superiority $k$; but if the largest magnitude does not exceed the sum of the other magnitudes, then the node is subequal. Then, whether we attend to distance or to magnitude, the three propositions hold good: (1) there cannot be more than one subequal node; (2) there cannot be more than two nodes each superior by unity; and if there is one such node, there is contiguous to it another such node; (3) starting from a node which is superior by more than unity, there is always a contiguous node which is of smaller superiority or else subequal; whence, as in the solution just referred to, there is always, as regards distance, a centre or a bicentre; and there is always, as regards magnitude, a centre or a bicentre.
\end{solution}

\newpage

\section*{On Mr Artemas Martin's First Question in Probabilities [Vol.~XXVII., pp.~89, 90]}

\begin{center}\textsc{By Professor Cayley.}\end{center}

The question was, ``A says that B says that a certain event took place; required the probability that the event did take place, $p_1$ and $p_2$ being A's and B's respective probabilities of speaking the truth.''

The solutions, referred to or given on pp.~77---79 of volume xxvii. of the Reprint, give the following values for the probability in question:
\[
\begin{array}{ll}
\text{Todhunter's Algebra} & \dfrac{p_1p_2}{p_1p_2 + (1-p_1)(1-p_2)}, \\[1.5ex]
\text{Artemas Martin} & \dfrac{p_1p_2}{p_1p_2 + (1-p_1)(1-p_2)}, \\[1.5ex]
\text{American Mathematicians and Woolhouse} & p_1p_2.
\end{array}
\]
It seems to me that the true answer cannot be expressed in terms of $p_1$ and $p_2$ only, but that it involves two other constants $\beta$ and $k$; and my value is
\[
\text{Cayley} \quad \frac{p_1p_2}{p_1p_2 + \beta(1-p_1)(1-p_2) + k(1-\beta)(1-p_1)}.
\]

In obtaining this I introduce, but I think of necessity, elements which Mr~Woolhouse calls extraneous and imperfect.

B told A that the event happened, or he did not tell A this; the only evidence is A's statement that B told him that the event happened; and the chances are $p_1$ and $1-p_1$. But, in the latter case, either B told A that the event did not happen, or he did not tell him at all; the chances (on the supposition of the incorrectness of A's statement) are $\beta$ and $1-\beta$; and the chances of the three cases are thus $p_1$, $\beta(1-p_1)$, and $(1-\beta)(1-p_1)$. On the supposition of the first and second cases respectively, the chances for the event having happened are $p_2$ and $1-p_2$; on the supposition of the third case (viz. here there is no information as to the event) the chance is the antecedent probability $k$, and the whole chance in favour of the event is
\[
p_1p_2 + \beta(1-p_1)(1-p_2) + k(1-\beta)(1-p_1).
\]
If $\beta = 1$, we have Todhunter's solution; if $\beta = 0$, and also $k=0$, we have the solution preferred by Woolhouse; but we do not (otherwise than by establishing between $k$ and $\beta$ a relation which is quite arbitrary) obtain Martin's solution.

The error in A says that B told him as to the event, and says further that the event did happen: the probability of the truth of this compound statement is $p_1$, where, in calling the probability $=p_1$, we mean that if A is speaking the truth, the event took place; this is to be regarded as a simple statement, and the probability of the truth of the statement to be taken to be $=p_1p_2$. Mr~Martin seems to introduce into his solution a hypothesis contradictory to the assumptions of the question.

I remark further that in my solution I assume that the event is of such a nature that, when there is any testimony in regard to it, the probability is determined by that testimony, irrespectively of the antecedent probability. This is quite consistent with the antecedent probability being, not zero, but $k$ (as it may very well be) indefinitely small; so if $k$ were $=0$, the whole probability $=p_1p_2$. But there is no absolutely reason for assigning any determinate value for $k$; and so the solutions which assume respectively $\beta=1$ and $\beta=0$, seem to me on this ground erroneous.

\newpage

\section*{Problems 4870--5625 [Vol.~XXIX., January to June, 1878]}

\subsection*{4870. (Proposed by Professor Cayley.)\---[Vol.~XXIX., p.~20.]}

If $\alpha, \beta, \gamma, \delta$; $\alpha_1, \beta_1, \gamma_1, \delta_1$ are such that
\[
(\alpha_1-\delta_1)(\beta_1-\gamma_1) = (\alpha-\delta)(\beta-\gamma), \quad
(\beta_1-\delta_1)(\gamma_1-\alpha_1) = (\beta-\delta)(\gamma-\alpha), \quad
(\gamma_1-\delta_1)(\alpha_1-\beta_1) = (\gamma-\delta)(\alpha-\beta);
\]
show that the three equations
\[
\frac{dx}{\sqrt{(x-\alpha)(x-\beta)(x-\gamma)(x-\delta)}} = \frac{dx_1}{\sqrt{(x_1-\alpha_1)(x_1-\beta_1)(x_1-\gamma_1)(x_1-\delta_1)}}
\]
such are equivalent relations between the constants, and show also that the equations in the foregoing differences $\alpha_1-\delta_1$, $\beta_1-\delta_1$, $\gamma_1-\delta_1$ may be so constituted with the particular integral of the differential equation.

\subsection*{5625. (Proposed by Professor Cayley.)\---[Vol.~XXIX., pp.~96, 97.]}

Given three conics passing through the four points; and on the first a point $A$, on the second a point $B$, and on the third a point $C$. It is required to find on the first a point $A'$, on the second a point $B'$, and on the third a point $C'$, such that the intersections of the lines $B'C'$ and $BC$, $C'A'$ and $CA$, $A'B'$ and $AB$ lie on the first conic; $C'A'$ and $CA$, $A'B'$ and $BA$, $B'C'$ and $BC$ lie on the second conic; and $A'B'$ and $CB$, $C'B'$ and $CA$ lie on the third conic.

\subsection*{5306. (Proposed by Professor Cayley.)\---[Vol.~XXXI., January to June, 1879, p.~38.]}

The equation
\[
\{(x+y+z)^2 - yz - zx - xy\}^2 = 4(2q+1)xyz(x+y+z)
\]
represents a trinodal quartic curve having the lines $x=0$, $y=0$, $z=0$, $x+y+z=0$ for its four bitangents; $X=0$, $Y=0$, $Z=0$ represent the sides of the triangle formed by the three nodes.

\subsection*{5387. (Proposed by Professor Cayley.)\---[Vol.~XXXII., July to December, 1879, p.~35.]}

Show that a cubic surface has at most 4 conical points, and a quartic surface at most 16 conical points.

\subsection*{5927. (Proposed by Professor Cayley.)\---[Vol.~XXXIII., January to July, 1880, p.~17.]}

If $\{a+\beta+\gamma+\ldots\}^n$ denote the expansion of $(a+\beta+\gamma+\ldots)^n$, retaining those terms only $Na^pb^qc^r\ldots$ in which $b+c+d+\ldots \geqq p-1$, $c+d+\ldots \geqq p-2$, \&c.\ \&c.; prove that
\[
a^n = (x+a)^n - \{a\}^1(x+a+\beta)^{n-1} + \{a+\beta\}^2(x+a+\beta+\gamma)^{n-2} - \{a+\beta+\gamma\}^3(x+a+\beta+\gamma+\delta)^{n-3} + \&\text{c.}\ldots(1).
\]

\newpage

\section*{Problems 6155--6800 [Vol.~XXXVI.--XXXVII., 1881--1882]}

\subsection*{6155. (Proposed by Professor Cayley.)\---[Vol.~XXXVI., 1881, p.~21.]}

Given, by means of their metrical coordinates, any two lines; it is required to find their inclination, shortest distance, and the coordinates of the line of shortest distance.

\textsc{N.B.}---If $\lambda, \mu, \nu$ are the inclinations of a line to three rectangular axes, and $\alpha, \beta, \gamma$ the coordinates to the same axes of a point on the line, then the metrical coordinates of the line are
\[
\begin{aligned}
a,\; b,\; c &= \cos\lambda,\; \cos\mu,\; \cos\nu, \\
f,\; g,\; h &= \beta\cos\nu - \gamma\cos\mu,\; \gamma\cos\lambda - \alpha\cos\nu,\; \alpha\cos\mu - \beta\cos\lambda,
\end{aligned}
\]
satisfying identically the relations
\[
a^2+b^2+c^2=1, \quad af+bg+ch=0.
\]

\subsection*{6470. (Proposed by Professor Cayley.)\---[Vol.~XXXVI., p.~64.]}

It is required, by a real linear transformation, to express the equation of a given cubic curve in the form
\[
\ldots
\]

\subsection*{6766. (Proposed by Professor Cayley.)\---[Vol.~XXXVI., pp.~106, 107.]}

Find the stationary and the double tangents of the curve $x^4+y^4+z^4=0$.

\begin{solution}[Solution by the Proposer.]
Take $l$, $m$ and $n$ fourth roots of $+1$; then the 28 double tangents are the lines $x=ly$, $x=lz$, $y=lz$, $(4+4+4=)12$ lines; and the lines $x+my+nz=0$, 16 lines; and the first 12 of these, each counted twice, are the 24 stationary tangents.

In fact, any one of the 12 lines is an osculating tangent, or line meeting the curve in 4 coincident points; it counts therefore once as a double tangent, and twice as a stationary tangent. There should consequently be 16 other double tangents; and it only needs to be shown that these are the 16 lines $x+my+nz=0$. Consider any one line $x+my+nz=0$; for its intersections with the curve $x^4+y^4+z^4=0$, we have
\[
(my+nz)^4 + y^4 + z^4 = 0,
\]
or, as this may be written,
\[
(my+nz)^4 + m^4y^4 + n^4z^4 = 0,
\]
viz. this is
\[
2(1,1,1\char`\,my,nz)^4 = 0,
\]
or, what is the same thing,
\[
2[(1,1,1\char`\,my,nz)^4] = 0;
\]
so that the line is a double tangent, the two points of contact being given by means of $(1,1,1\char`\,my,nz)^2 = 0$; viz. $\omega$ being an imaginary cube root of unity, we have $nz = \omega my$ or $\omega^2 my$; and thence
\[
x : y : z = -m : 1 : -\frac{n}{\omega m} \quad\text{or}\quad -m : 1 : -\frac{n}{\omega^2 m},
\]
for the points of contact, for
\[
x+my+nz=0, \quad x^4+y^4+z^4=0.
\]
\end{solution}

\newpage

\subsection*{6800. (Proposed by W.~J.~C.~Miller, B.A.)\---[Vol.~XXXVII., 1882, p.~74.]}

Prove that, if
\[
\frac{ayz}{y^2+z^2} = \frac{bzx}{z^2+x^2} = \frac{cxy}{x^2+y^2},
\]
then
\[
a^2+b^2+c^2 = abc+4.
\]

\begin{note}[By Professor Cayley. Note on Question 6800.]
The identity given by the solution is a very interesting one. Writing $(a,b,c,d)$ as the coordinates of a point in space, and $(x,y,z)$ for a given system of values of $(x,y,z)$; and considering $(a,b,c,d)$ as coordinates of a point in the plane, it may be shown without difficulty that, to a given point $(x,y,z)$, there correspond two systems of values of $(a,b,c,d)$ satisfying
\[
4d^2 - d(a^2+b^2+c^2) + abc = 0;
\]
to a given system of values $(a,b,c,d)$ there corresponds, it is clear, a single point $(x,y,z)$ in the plane. We have thus a correspondence between the points of the cubic surface $4d^2 - d(a^2+b^2+c^2) + abc = 0$, and the points of the plane; the plane and cubic surface have thus a $(1,2)$ correspondence with each other.
\end{note}

\newpage

\section*{Problems 5244--5689 [Vol.~XXXVIII.--XXXIX., 1883]}

\subsection*{5244. (Proposed by Professor Cayley.)\---[Vol.~XXXVIII., 1883, pp.~87--89.]}

Writing for shortness
\[
\begin{aligned}
A &= a^2+\beta^2+\gamma^2+\delta^2, \\
G &= \beta^2+\gamma^2+\delta^2-a^2, \\
H &= \gamma^2+\delta^2+a^2-\beta^2, \\
F &= \delta^2+a^2+\beta^2-\gamma^2, \\
L &= a^2+\beta^2-\gamma^2-\delta^2, \\
M &= a^2-\beta^2+\gamma^2-\delta^2, \\
N &= a^2-\beta^2-\gamma^2+\delta^2,
\end{aligned}
\]
show that the equation
\[
\begin{aligned}
&LMN(x^2+y^2+z^2+w^2)^2 \\
&\quad - MN(F^2+2L)(y^2z^2+x^2w^2) + NL(G^2+2M)(z^2x^2+y^2w^2) \\
&\quad + LM(H^2+2N)(x^2y^2+z^2w^2) - 2a^2\beta^2\gamma^2\delta^2\, FGHA\, xyzw = 0
\end{aligned}
\]
belongs to a 16-nodal quartic surface, having the nodes
\[
\begin{array}{c|cccc}
& x & y & z & w \\ \hline
1 & a & \beta & \gamma & \delta \\
2 & a & \beta & \gamma & -\delta \\
3 & a & \beta & -\gamma & \delta \\
4 & a & \beta & -\gamma & -\delta \\
5 & a & -\beta & \gamma & \delta \\
6 & a & -\beta & \gamma & -\delta \\
7 & a & -\beta & -\gamma & \delta \\
8 & a & -\beta & -\gamma & -\delta \\
9 & -a & \beta & \gamma & \delta \\
10 & -a & \beta & \gamma & -\delta \\
11 & -a & \beta & -\gamma & \delta \\
12 & -a & \beta & -\gamma & -\delta \\
13 & -a & -\beta & \gamma & \delta \\
14 & -a & -\beta & \gamma & -\delta \\
15 & -a & -\beta & -\gamma & \delta \\
16 & -a & -\beta & -\gamma & -\delta
\end{array}
\]
and the sixteen singular tangent planes represented by the equations
\[
(a,\beta,\gamma,\delta)(x,y,z,w) = 0, \quad \&\text{c.}
\]

\subsection*{7190. (Proposed by Professor Wolstenholme, M.A.)\---[Vol.~XXXVIII., pp.~87--89.]}

If $x,y,z$ be three quantities satisfying the two symmetrical equations
\[
yz+zx+xy=0, \quad x^2+y^2+z^2+3xyz=0,
\]
prove that (1) they will satisfy also one of the two pairs of semi-symmetrical expressions
\[
yz^2+z^2x+xy^2 = (y-z)(z-x)(x-y) = +xyz,
\]
\[
yz^2+z^2x+xy^2 = (y-z)(z-x)(x-y) = -xyz;
\]
and (2) one set of the following equations will also be satisfied:
\[
\begin{aligned}
x^2+yz-z^2&=0, & y^2+zx-x^2&=0, & z^2+xy-y^2&=0; \\
x^2+yz-y^2&=0, & y^2+zx-z^2&=0, & z^2+xy-x^2&=0.
\end{aligned}
\]

\begin{solution}[Solution by Professor Cayley.]
The two symmetrical equations represent a conic and a cubic respectively; and if we denote by $\alpha$ a root of the equation
\[
u^3+\nu^2-2\nu-1=0,
\]
then the other two roots of this equation are $\beta, \gamma$; viz. if $\alpha^3+\alpha^2-2\alpha-1=0$, then we have
\[
(u-\alpha)(u+1+\tfrac{1}{\alpha})(u+\tfrac{1}{\alpha}-\alpha) = u^3+u^2-2u-1,
\]
an identity which is easily verified.

It then may be remarked that, if the last left-hand side of the six mentioned equation which remains unaltered when $\alpha$ is changed into $\frac{1}{\alpha}$ or $-\frac{1}{\alpha}-1$ is
\[
(u-\alpha)(u-\phi\alpha)(u-\psi\alpha),
\]
then the coordinates can be expressed indifferently in terms of any one of the roots $(\alpha,\beta,\gamma)$, viz. the coordinates are
\[
(\alpha^2-1, \; -\alpha, \; -1), \quad (-1, \; \alpha^2-1, \; -\alpha), \quad (-\alpha, \; -1, \; \alpha^2-1), \quad \ldots\; (1,2,3),
\]
or they are
\[
(\alpha^2-1, \; -1, \; -\alpha), \quad (-\alpha, \; \alpha^2-1, \; -1), \quad (-1, \; -\alpha, \; \alpha^2-1), \quad \ldots\; (4,5,6);
\]
the coordinates of the points equal to the like expressions in $\beta$ and $\gamma$ respectively, as shown by the attached numbers.

Thus, writing $x,y,z = \alpha^2-1, \; -\alpha, \; -1$, we find
\[
yz+zx+xy = -\alpha - \alpha^2 + 1 - \alpha^2 + \alpha = -(\alpha^3+\alpha^2-2\alpha-1) = 0,
\]
\[
\begin{aligned}
x^2+y^2+z^2+4xyz &= (\alpha^2-1)^2 + \alpha^2 + 1 + 4\alpha(\alpha^2-1) \\
&= \alpha^6 - 3\alpha^4 + 3\alpha^2 + 3\alpha - 4\alpha - 2 = (\alpha^3+\alpha^2-2\alpha-1)(\alpha^3-\alpha^2+2) = 0,
\end{aligned}
\]
which verifies the formulas for the six points of intersection.

Take, again, $x,y,z = \alpha^2-1, \; -\alpha, \; -1$; then we find
\[
yz^2+z^2x+xy^2 = -\alpha - (\alpha^2-1)^2 + \alpha^2(\alpha^2-1) = \alpha^4 - \alpha - 1,
\]
\[
y^2z+z^2x+x^2y = +(\alpha^2-1) - \alpha(\alpha^2-1)^2 = -\alpha^5 + 2\alpha^3 - \alpha - 1.
\]
Or, since $\alpha^3 = -\alpha^2+2\alpha+1$, and thence $\alpha^4 = 3\alpha^2-\alpha-1$, $\alpha^5 = -4\alpha^2+5\alpha+3$, the last equation becomes
\[
y^2z+z^2x+x^2y = 2\alpha^2 - 2\alpha - 2.
\]
We have also
\[
xyz = \alpha^3 - \alpha = -\alpha^2 + \alpha + 1;
\]
hence the point in question is situate on each of the cubics
\[
yz^2+z^2x+xy^2+xyz=0, \quad y^2z+z^2x+x^2y+2xyz=0,
\]
and
\[
y^2z+z^2x+x^2y-2(yz^2+z^2x+xy^2)=0;
\]
and this, of course, shows that each of the three cubics; and the points 4, 5, 6 are all of them situate on each of the three cubics $yz^2+z^2x+xy^2+2xyz=0$, $y^2z+z^2x+x^2y+xyz=0$, $y^2z+z^2x+x^2y-2(yz^2+z^2x+xy^2)=0$.

Again, from the values $x,y,z = \alpha^2-1, \; -\alpha, \; -1$, we have
\[
x^2+yz-y^2=0, \quad z^2+xy-x^2=0, \quad y^2+zx-z^2=0;
\]
on each of these conics; similarly the point 2 lies on each of the same conics; and the point 3 lies on each of the same conics that is, the conics in question have in common the points 1, 2, 3.

In like manner, the conics $x^2+yz-z^2=0$, $y^2+zx-x^2=0$, $z^2+xy-y^2=0$ have in common the points 4, 5, 6.

The general result is that the given conic and the cubic meet in six points forming two groups of points $(1,2,3)$ and $(4,5,6)$; through the points $(1,2,3)$ we have three cubics and three conics; and through the points $(4,5,6)$ we have three cubics and three conics.

If in the equation $u^3+u^2-2u-1=0$, whose roots are $(\alpha), \beta, \gamma$, we put $u=2\cos\theta$, the equation becomes
\[
2(3\cos\theta+\cos3\theta) + 2(1+\cos2\theta) - 4\cos\theta - 2 = 0,
\]
or
\[
2\cos3\theta + 2\cos2\theta + 2\cos\theta = 0;
\]
or $\dfrac{\sin\frac{3\theta}{2}}{\sin\frac{\theta}{2}} = 0$; or the three roots are $2\cos\frac{2}{7}\pi$, $2\cos\frac{4}{7}\pi$, $2\cos\frac{6}{7}\pi$.

The two equations $yz+zx+xy=0$, $x^2+y^2+z^2+3xyz=0$, are satisfied if $x:y:z = \text{these three roots in any order}$, giving the six solutions.

The semi-symmetrical systems are satisfied, the one by $x:y:z = \cos\frac{2}{7}\pi : \cos\frac{4}{7}\pi : \cos\frac{6}{7}\pi$, or $y:z:x$, or $z:x:y$; and the other by $x:y:z = \cos\frac{2}{7}\pi : \cos\frac{6}{7}\pi : \cos\frac{4}{7}\pi$, or $y:z:x$, or $z:x:y$, or $x:z:y$, or $y:x:z$, or $z:y:x$.
\end{solution}

\newpage

\subsection*{5689. (Proposed by Professor Cayley.)\---[Vol.~XXXIX., 1883, p.~31.]}

Show that (1) the apparent contour of a Steiner's surface $(2xyz+y^2z^2+z^2x^2+x^2y^2=0)$, as seen from an exterior point on a nodal line (say the axis of $z$), projected on the plane of the other two nodal lines, is an ellipse passing through the four points $(\pm1,0)$ and $(0,\pm1)$; and (2) find the surface-contour, or curve of contact, of the cone and surface.

\subsection*{4722. (Proposed by Professor Cayley.)\---[Vol.~XXXIX., p.~49.]}

Show that
\begin{enumerate}[label=\arabic*.]
\item the conditions in regard to the reality of the roots of the equation
\[
(x^2-a)^2 + 16A(x-m) = 0,
\]
are, if
\[
(4m^2-3a)^2 - (8m^2-9ma-27A)^2 = -,
\]
then the roots are two real, two imaginary; but if
\[
(4m^2-3a)^2 - (8m^2-9ma-27A)^2 = +,
\]
then, if simultaneously $A=+$ and $A(ma-9A)=+$, the roots are all real, but otherwise they are imaginary.

\item If the roots of the foregoing equation are all imaginary, then for any value whatever of $y$, the roots of the equation
\[
(x^2+y^2-a)^2 + 16A(x-my) = 0
\]
are all imaginary.
\end{enumerate}

Using the term ``Cassinian'' to denote a bi-circular quartic having four foci in a right line; show that the equation of a Cassinian having its four foci at the points $x=a$, $x=b$, $x=c$, $x=d$ on the axis of $X$, may be written in the four equivalent forms
\[
\begin{aligned}
&\tau(d-c)\cdot A_1 + \rho(b-d)\cdot B_1 + \sigma(c-b)\cdot C_1 + \tau(d-c)\cdot D_1 = 0, \\
&\tau(c-d)\cdot A_1 + \rho(d-a)\cdot C_1 + \sigma(a-c)\cdot D_1 = 0, \quad \&\text{c.}, \&\text{c.},
\end{aligned}
\]
where $A_1$, $B_1$, $C_1$, $D_1$ are the distances from the four foci and $\rho$, $\sigma$, $\tau$ are connected by the equation
\[
\rho^2(a-d)(b-c) + \sigma^2(b-d)(c-a) + \tau^2(c-d)(a-b) = 0.
\]
Show also that the curve has, at right angles to the axis of $x$, two double tangents, the equation whereof is any one of the three equivalent forms
\[
(a+d-2x)(b+c-2x) : (b+d-2x)(c+a-2x) : (c+d-2x)(a+b-2x) = \rho^2 : \sigma^2 : \tau^2.
\]

\newpage

\section*{Problems 7376--5421 [Vol.~XL.--XLI., 1884]}

\subsection*{7376. (Proposed by Professor Cayley.)\---[Vol.~XL., 1884, p.~32.]}

Show how the construction of a regular heptagon may be made to depend on the trisection of the angle $\cos^{-1}\left(\dfrac{\sqrt{7}-1}{2\sqrt{2}}\right)$.

\subsection*{7352. (Proposed by Professor Cayley.)\---[Vol.~XL., p.~110.]}

Denoting by $x,y,z,\xi,\eta,\zeta$ homogeneous linear functions of four coordinates, such that identically
\[
x+y+z+\xi+\eta+\zeta = 0,
\]
where $af=bg=ch=1$; show that
\[
\sqrt{(x\xi)} + \sqrt{(y\eta)} + \sqrt{(z\zeta)} = 0
\]
is the equation of a quartic surface having the sixteen singular tangent planes (each touching it along a conic)
\[
\begin{aligned}
&ax+by+cz+f\xi+g\eta+h\zeta=0, \quad x+y+z=0, \quad x+\eta+\zeta=0, \\
&ax+by+cz=0, \quad \xi+y+\zeta=0, \quad x+y+\xi=0, \quad ax+g\eta+cz=0, \\
&x=0, \quad y=0, \quad z=0, \quad \xi=0, \quad \eta=0, \quad \zeta=0, \\
&\frac{1-bc}{1-ca}, \quad \frac{1-ca}{1-ab}, \quad \frac{1-ab}{1-fg}, \quad \frac{1-gh}{1-hf}, \quad \frac{1-hf}{1-fg}.
\end{aligned}
\]

\subsection*{5421. (Proposed by Professor Cayley.)\---[Vol.~XLI., 1884, p.~37.]}

Suppose
\[
\delta x = m_1(x-a_1), \; m_2(x-a_2), \; m_3(x-a_3), \; m_4(x-a_4),
\]
where, for any given value of $x$, we write $+$, $-$ or $0$, according as the linear function is positive, negative, or zero; and where the order of the terms is not attended to.

If $x$ is any one of the values $a_1, a_2, a_3, a_4$, the corresponding $\delta$ is $+++-$, $++-+$, $+-++$, $-+++$; and if $I$ denote indifferently the first or the second form, and $R$ denote indifferently the third or the fourth form, then it is to be shown that the four $\delta$'s are $R,R,R,R$ or else $R,R,I,I$.

\newpage

\section*{Problems 8340--5271 [Vol.~XLIV.--XLVII., 1886--1887]}

\subsection*{8340. (Proposed by F.~Morley, B.A.)\---[Vol.~XLIV., 1886, p.~109.]}

Show that (1) on a chess-board the number of squares visible is 204, and the number of rectangles (including squares) visible is 1,296; and (2) on a similar board, with $n$ squares in each side, the number of visible squares is the sum of the first $n$ square numbers, and the number of rectangles (including squares) is the sum of the first $n$ cube numbers.

\begin{solution}[Solution by Professor Cayley.]
In a board of $n^2$ squares, the number of squares of breadth $n-r+1$ is $=r$; from each other of pairs of vertical lines at a distance $n-r+1$, and the number of pairs of horizontal lines at a distance $n-s+1$ is $=s$. Hence the number of rectangles of breadth $n-r+1$ and depth $n-s+1$, or say the number of $(n-r+1)(n-s+1)=rs$.

For instance, $n=4$, the number of rectangles $4^2$, $4\cdot3$, $3\cdot4$, \&c., is shown in the diagram; hence the whole number of rectangles is $(1+2+3+4)^2 = 1^3+2^3+3^3+4^3$, and so for any value of $n$.

The same diagram shows that the whole number is $=1^2+2^2+3^2+4^2$ and so for any value of $n$.
\end{solution}

\subsection*{8636. (Proposed by Mahendra Nath Ray, M.A., LL.B.)\---[Vol.~XLVII., 1887, pp.~49, 50.]}

Show that the following equations are satisfied by the same value of $x$, and find this value:
\[
\begin{aligned}
ax(x^2-a^2)^{\frac{1}{2}} + bx(x^2-b^2)^{\frac{1}{2}} + cx(x^2-c^2)^{\frac{1}{2}} &= 2abc, \\
2(x^2-a^2)^{\frac{1}{2}}(x^2-b^2)^{\frac{1}{2}}(x^2-c^2)^{\frac{1}{2}} &= x(a^2+b^2+c^2-2x^2).
\end{aligned}
\]

\begin{solution}[Solution by Professor Cayley.]
The second equation rationalised gives
\[
4(x^2-a^2)(x^2-b^2)(x^2-c^2) = x^2(a^2+b^2+c^2-2x^2)^2,
\]
that is,
\[
4x^6 - 4x^4(a^2+b^2+c^2) + 4x^2(b^2c^2+c^2a^2+a^2b^2) - 4a^2b^2c^2 = 4x^6 - 4x^4(a^2+b^2+c^2) + x^2(a^2+b^2+c^2)^2;
\]
if, for shortness,
\[
V = -a^4-b^4-c^4+2b^2c^2+2c^2a^2+2a^2b^2.
\]
We thence find
\[
\begin{aligned}
V(x^2-a^2) &= a^2(-a^2+b^2+c^2)^2, \\
V(x^2-b^2) &= b^2(a^2-b^2+c^2)^2, \\
V(x^2-c^2) &= c^2(a^2+b^2-c^2)^2;
\end{aligned}
\]
and therefore
\[
V\sqrt{a^2b^2c^2(x^2-a^2)} = 2a^2b^2c(-a^2+b^2+c^2), \quad \&\text{c.}
\]
Or, assuming the sign of the square roots,
\[
\begin{aligned}
V\sqrt{ax(x^2-a^2)} &= 2abc(-a^2+b^2+c^2)^{\frac{1}{2}}, \\
V\sqrt{bx(x^2-b^2)} &= 2abc(+a^2b^2-b^4+b^2c^2)^{\frac{1}{2}}, \\
V\sqrt{cx(x^2-c^2)} &= 2abc(c^2a^2+b^2c^2-c^4)^{\frac{1}{2}};
\end{aligned}
\]
whence, adding, the whole divides by $V$ and we have
\[
ax(x^2-a^2)^{\frac{1}{2}} + bx(x^2-b^2)^{\frac{1}{2}} + cx(x^2-c^2)^{\frac{1}{2}} = 2abc.
\]
Observe that the second equation rationalised gives an equation of the form $(x^2,1)^4=0$, that is, the foregoing value of $x^2$ is one of the four values of $x^2$.
\end{solution}

\newpage

\subsection*{5271. (Proposed by Professor Cayley.)\---[Vol.~XLVII., 1887, p.~141.]}

If $\omega$ be an imaginary cube root of unity, show that, if
\[
\frac{(1-\omega^2)x + \omega^2(1-\omega x^2)}{(1-\omega^2) - \omega(1-\omega x^2)} \cdot \frac{dy}{dx} = \frac{(1-\omega)(1+\omega x^2)}{(1-y^2)(1+\omega y^2)},
\]
then
\[
\frac{dy}{\sqrt{(1-y^2)(1+\omega y^2)}} = \frac{dx}{\sqrt{(1-x^2)(1+\omega x^2)}};
\]
and explain the general theory.

\newpage

\section*{Problem 3105 [Vol.~L., 1889, p.~189]}

\subsection*{3105. (Proposed by Professor Cayley.)\---[Vol.~L., 1889, p.~189.]}

The following singular problem of partitions arises out of the geometrical theory given in Cremona's Memoir, ``Sulle trasformazioni geometriche delle figure piane,'' Mem.~di Bologna, tom.~v.~(1865).

A number is to be partitioned on the same type $1\cdot11 + 8\cdot2$ as follows: for instance, $27 = 1\cdot11 + 8\cdot2$.

(Observe that the partitionment is to be made up in any manner with the literal factors $a,b,c,d,e,f,g,h,i$; always the positive integers as the product of 27.) Forming, then, symmetrical products, we have $a^9(bcdefghi)$.

But, to take another example, $37 = 0\cdot11 + 3\cdot8 + 1\cdot5 + 4\cdot2 = 1\cdot11 + 0\cdot8 + 4\cdot5 + 3\cdot2$.

The first of these (symmetrically as above) gives on its own type, and the second gives on its own type; though it may be on the second type; viz. the partitions of the two products respectively are:

First product on second type: $abcdefgh$, $ab$, $ac$, $ad$, $ae$, $af$, $ag$, $ah$, $ai$, $bc$, $bd$, $be$, $bf$, $bg$, $bh$, $bi$.

Second product on first type: $a^2bcdegh$, $a^2bcdefh$, $a^2bcdefg$, $abcde$, $abcdf$, $abcdg$, $abcdh$, $abcde$, $abcdf$, $abcdg$, $abcdh$.

so that in the first example the type is sibi-reciprocal, but in the second example there are two conjugate types. Other examples are:

\begin{center}
\begin{tabular}{c|ccccc}
Parts & 55 & 53 & 48 & 64 & 55 \\ \hline
& 2 & 14 & 3 & 1 & 3 \\
& 5 & 2 & 3 & 6 & 8 \\
& 3 & 2 & 6 & 5 & 7 \\
& 1 & 5 & 2 & 3 & 2 \\
& 56 & 2 & 6 & 5 & 7 \\
& 1 & 5 & 2 & 3 & 2 \\ \hline
No. & 1 & 3 & 4 & 1 & 1 \\
& 13 & 1 & 17 & 1 & 20
\end{tabular}
\end{center}

viz. the first four columns give each of them a sibi-reciprocal type, but the last two double columns give conjugate types. It is required to investigate the general solution.

\newpage

\section*{Problems 3304--3481 [Vol.~L., 1889, pp.~191, 192]}

\subsection*{3304. (Proposed by Professor Cayley.)\---[Vol.~L., p.~191.]}

The coordinates $x,y,z$ being proportional to the perpendicular distances from the sides of an equilateral triangle, it is required to trace the curve
\[
(y-z)\sqrt{x} + (z-x)\sqrt{y} + (x-y)\sqrt{z} = 0.
\]

[Prof.~Cayley remarks that the curve which presents itself in the following theorem, communicated to him (with a demonstration) several years ago by Mr~J.~Griffiths, is a particular case of a more general curve: The locus of a point $(x,y,z)$ such that its pedal circle (that is, the circle which passes through the feet of the perpendiculars drawn from the point in question upon the sides of the triangle of reference) touches the nine-point circle, is the sextic curve
\[
\begin{aligned}
&y^2z^2\cos^2A(x\cos A)^2 + z^2x^2\cos^2B(y\cos B)^2 + x^2y^2\cos^2C(z\cos C)^2 \\
&\quad - 2yz\cos A(x\cos A)(y\cos B)(z\cos C) - \cdots = 0.
\end{aligned}
\]
To trace this more general curve would be an interesting problem.]

\subsection*{3481. (Proposed by Professor Cayley.)\---[Vol.~L., p.~192.]}

Find, in the Hamiltonian form, the equations for the motion of a particle acted on by a central force in a plane.

\newpage

\section*{Problems 10716--3186 [Vol.~LXI., 1894, pp.~122--124]}

\subsection*{10716. (Proposed by Professor Cayley.)\---[Vol.~LXI., 1894, pp.~122, 123.]}

In a hexahedron $ABCDA'B'C'D'$ the edges $AA'$, $BB'$, $CC'$, $DD'$ intersect in four points, say $AA'$, $DD'$ in $\alpha$; $BB'$, $CC'$ in $\beta$; $CC'$, $DD'$ in $\gamma$; $AA'$, $BB'$ in $\delta$: that is, starting with the duad of lines $\alpha\beta$, $\gamma\delta$, which join the extremities of the edges, the four edges are the lines joining the extremities of a duad. The question arises, ``Given two tetrads $ABCD$, $A'B'C'D'$, when are the lines joining the eight of the twelve edges of a hexahedron?'' Given the tetrahedron $\alpha\beta\gamma\delta$ and the two tetrads of joining lines, so that the duads $\alpha\beta$, $\gamma\delta$; $\alpha\gamma$, $\beta\delta$; $\alpha\delta$, $\beta\gamma$ which join the extremities of the edges are the four duads; and the lines $AB$, $CD$, $A'B'$, $C'D'$ are the edges $AD$, $BC$, $A'D'$, $B'C'$.

The duad $\alpha\beta$, $\gamma\delta$ is considered which determines the relative position of the two finite lines $\alpha\beta$ and $\gamma\delta$.

\subsection*{3162. (Proposed by Professor Cayley.)\---[Vol.~LXI., p.~123.]}

By a proper determination of the coordinates, the skew cubic through any six given points may be taken to be
\[
x:y:z:w = 1:t:t^2:t^3,
\]
or, what is the same thing, the coordinates of the six points may be taken to be $(1,t_i,t_i^2,t_i^3)$, $i=1,2,\ldots,6$. Assuming this, it is required to show that if
\[
V = 9xyzw - 4xz^3 - 4yw^3 + 3z^2w^2 - x^2w^2,
\]
then the equation of the Jacobian surface of the six points is
\[
\begin{aligned}
&(2xp_2-yp_1-2w)\partial_xV \\
&\quad + (3xp_3-2yp_2-zp_1+w)\partial_yV \\
&\quad + (2xp_4-yp_3-2zp_2+wp_1)\partial_zV \\
&\quad + (xp_4-2yp_3+zp_2+2w)\partial_wV = 0,
\end{aligned}
\]
where $p_1=\sum t_i$, $p_2=\sum t_it_j$, \&c.

\subsection*{3186. (Proposed by Professor Cayley.)\---[Vol.~LXI., p.~124.]}

An unclosed polygon of $(m+1)$ vertices is constructed as follows: the abscissae of the several vertices are $0,1,2,\ldots,m$; and to the abscissa $k$ the ordinate is the chance of $m+k$ heads in $2m$ tosses of a coin; then the areas of any very large and small portions of the polygon, what information is afforded in regard to the general results of the Theory of Probabilities?

\subsection*{3229. (Proposed by Professor Cayley.)\---[Vol.~LXI., p.~124.]}

It is required to find the value of the elliptic integral $F(c,\theta)$ when $c$ is very nearly $=1$ and $\theta$ very nearly $=\frac{1}{2}\pi$; that is, the value of
\[
\int_0^{\frac{\pi}{2}-b} \frac{d\theta}{\sqrt{\sin^2(a)+\cos^2(a)\sin^2(\theta)}}
\]
where $a,b$ are each of them indefinitely small.

\textsc{N.B.}---Observe that, for $a=0$, $b$ small, the value is $=\log\cot\frac{1}{2}b$; and for $b=0$, $a$ small, the value is equal $\log\frac{4}{a}$, and for $a=0$, $b=0$, the value is $=\infty$.

\newpage

\section*{Index of Problems [pp.~615--616]}

\begin{center}
\textsc{In the following Contents, the Problems are referred to each by its number and the page; an asterisk shews that there is no number; a dash shews that no solution is given.}
\end{center}

\medskip

\begin{center}
\begin{tabular}{r@{\quad}l@{\quad}r@{\quad}l}
\multicolumn{4}{c}{\textsc{Page}} \\[0.5ex]
3002 & Collins & 566 & Systems of circles \\
3144 & Cayley & 568 & Locus in piano \\
3120 & \ditto & 569 & Jacobian of quadric surfaces \\
3249 & \ditto & 569 & Sets of four points on a conic \\
3206 & \ditto & 570 & Arrangements of points and lines \\
3278 & \ditto & 570 & Arrangements of numbers \\
3329 & \ditto & \ditto & Substitutions and permutations \\
3356 & \ditto & 574 & Roots of quartic equation \\
3507 & \ditto & 575 & Quadric cones through given points \\
3536 & \ditto & \ditto & Particle under central forces \\
3591 & \ditto & \ditto & Relation between areas of triangles \\
2652 & \ditto & \ditto & Parallel surfaces of an ellipsoid \\
3677 & \ditto & 576 & Angle between normal and bisector of chord \\
3564 & \ditto & \ditto & Minimum circle enclosing three points \\
3875 & \ditto & \ditto & Description of curves \\
3430 & \ditto & 578 & Negative pedals of ellipse \\
4298 & Glaisher & \ditto & Quadrilateral inscribable in circle \\
4392 & Roberts & \ditto & Symmetrical determinant \\
4354 & Tucker & 581 & Geometrical interpretation (note) \\
4458 & Cayley & 582 & Two quartic curves \\
4520 & Miller & \ditto & \ditto \\
4620 & Evans & 586 & Use of Degen's tables \\
4528 & Cayley & 587 & Expectation \\
4581 & Wilkinson & 588 & A question in chances \\
4638 & Cayley & 589 & Envelope of family of quadrics \\
4694 & \ditto & \ditto & Resultant of forces \\
4793 & Wolstenholme & 590 & Relation among derivatives of a function \\
4752 & Cayley & 592 & Cubic curve \\
4946 & \ditto & 594 & Attraction of lens-shaped body \\
5020 & Woolhouse & \ditto & Algebraical theorem \\
5079 & Cayley & 596 & Bicircular quartic \\
5130 & \ditto & 597 & Bicircular quartic \\
5208 & Sylvester & 598 & Trees \\
5306 & Cayley & 600 & A question in probabilities \\
4870 & \ditto & 601 & A system of equations \\
5625 & \ditto & 602 & Conics and lines \\
5387 & \ditto & \ditto & Conical points \\
5927 & \ditto & 603 & Algebraical theorem \\
6155 & \ditto & \ditto & Coordinates of a line \\
6470 & \ditto & \ditto & Equation of a cubic curve \\
6766 & Miller & \ditto & Singular tangents of a quartic \\
6800 & Cayley & 604 & Geometrical interpretation (note) \\
5244 & \ditto & \ditto & 16 nodal quartic surface \\
7190 & Wolstenholme & 605 & A Steiner's surface \\
5689 & Cayley & 607 & Reality of roots of a quartic \\
4722 & \ditto & 608 & Equation of a Cassinian \\
4387 & \ditto & 609 & Construction of a heptagon \\
7376 & \ditto & \ditto & Equation of a quartic surface \\
7352 & \ditto & \ditto & Algebraical theorem \\
5421 & \ditto & \ditto & \ditto \\
8340 & Morley & \ditto & Topology of chess-board \\
8636 & Nath Ray & 610 & Solution of equations \\
5271 & Cayley & 611 & Elliptic-function transformation \\
3105 & \ditto & \ditto & Partitions \\
3304 & \ditto & 612 & Curve of sixth order \\
3481 & \ditto & 613 & Hamiltonian equations of central orbit \\
10716 & \ditto & \ditto & Edges of a hexahedron \\
3162 & \ditto & \ditto & Jacobian surface of six points \\
3186 & \ditto & \ditto & Probability \\
3229 & \ditto & \ditto & Elliptic integral
\end{tabular}
\end{center}

\vfill

\begin{center}
\textsc{End of Vol.~X}

\medskip

CAMBRIDGE: PRINTED BY J.~AND C.~F.~CLAY, AT THE UNIVERSITY PRESS.
\end{center}

\end{document}
